\section{Security, Supply Chain, and Permissions}

Agentic coding expands the write surface. The security posture must therefore be structural. CISA/NSA guidance pushes memory-safe languages as a way to reduce memory-related vulnerability classes \cite{cisaMemorySafeLanguages2025}. Veracode's GenAI security results show syntax success does not imply secure code \cite{veracodeGenaiCodeSecurity2025}. GitGuardian's secret-sprawl evidence shows AI-churned repositories need stronger secret gates, not more trust \cite{gitguardianSecretsSprawl2026}.

\subsection{Security Considerations}

Jankurai evidence can itself become sensitive. A conformant implementation MUST treat logs, screenshots, traces, prompts, browser storage, dependency reports, and security scan output as evidence with privacy rules, not as disposable CI noise. Evidence bundles should avoid secrets by construction, redact customer data, hash large binary artifacts, separate public and restricted receipts, carry retention metadata, and fail release when required redaction status is missing.

\begin{table}[t]
\caption{Evidence privacy classes.}
\label{tab:evidence-privacy}
\centering
\scriptsize
\begin{tabularx}{\columnwidth}{L{0.22\columnwidth} Y}
\toprule
\JKTableHeader
\textbf{Class} & \textbf{Merge and publication rule} \\
\midrule
Public & Safe to publish with report or badge. \\
Internal & May satisfy CI evidence but not public registry evidence. \\
Restricted & Requires access controls, retention bounds, and redacted public pointer. \\
Secret-bearing & Fails the lane until removed, rotated where needed, and replaced by clean evidence. \\
\bottomrule
\end{tabularx}
\end{table}

\begin{table}[t]
\caption{Security obligations by conformance level.}
\label{tab:security-hl-obligations}
\centering
\scriptsize
\begin{tabularx}{\columnwidth}{L{0.16\columnwidth} Y}
\toprule
\JKTableHeader
\textbf{Level} & \textbf{Security obligation} \\
\midrule
HL1 & Advisory scan output records missing security tools. \\
HL2 & Secret sprawl, generated mutation, destructive migration, and overbroad agency can block. \\
HL3 & High/critical security findings require current receipts and witness binding. \\
HL4 & Security posture cannot regress without waiver and repair queue. \\
HL5 & Release evidence includes security, provenance, redaction, and rollback readiness. \\
\bottomrule
\end{tabularx}
\end{table}

\begin{table*}[t]
\caption{Permission receipt matrix.}
\label{tab:permission-matrix}
\centering
\small
\begin{tabularx}{\textwidth}{L{0.22\textwidth} Y Y}
\toprule
\JKTableHeader
\textbf{Permission surface} & \textbf{Allowed evidence} & \textbf{Blocking concern} \\
\midrule
Repo reads & Changed-path inventory, source-trust labels, context-pack receipt. & Reading secrets, private transcripts, or unrelated customer data into context. \\
Tracked writes & Owner route, changed paths, test-map lane, proof receipt. & Writes outside declared scope. \\
Generated zones & Source contract path, generator command, artifact digest. & Hand-edited generated output. \\
Dependency install & Lockfile diff, rationale, SCA/OSV result, SBOM/provenance. & Unreviewed package drift or install script risk. \\
Network & Destination, purpose, lane, timeout, cache policy. & Unbounded exfiltration or live-service dependence in proof. \\
Secrets & Named secret broker, redaction proof, no raw value in artifacts. & Secret material in prompts, logs, screenshots, or fixtures. \\
Migrations & Migration analysis, rollback/backfill/lock proof, DB lane receipt. & Destructive migration without safety evidence. \\
Privileged browser automation & Route, account class, storage isolation, screenshot redaction. & Capturing private state or bypassing auth controls. \\
Push/PR & Branch name, diff scope, proof summary, reviewer route. & Publishing unreviewed or out-of-scope changes. \\
Auto-merge & Out of scope for this edition. & Merge without current witness and human/CI policy gate. \\
\bottomrule
\end{tabularx}
\end{table*}

\begin{table*}[t]
\caption{Security controls for agent-native repositories.}
\label{tab:security-controls}
\centering
\small
\begin{tabularx}{\textwidth}{L{0.18\textwidth} Y Y}
\toprule
\JKTableHeader
\textbf{Risk} & \textbf{Control} & \textbf{Repair evidence} \\
\midrule
Prompt injection and hostile context & Source hierarchy: root instructions, local owner rules, trusted docs, then untrusted issue/web text. & Receipt names ignored instruction, trusted rule, and resulting safe action. \\
Secrets and credentials & Secret scanning, no secrets in UI or docs, redacted logs, no model prompt leakage. & Scanner artifact, redaction test, finding closure. \\
Dependency sprawl & Lockfiles, SCA, SBOM/provenance, dependency rationale for new packages. & Package diff, risk reason, scanner result \cite{githubSecretScanning2026,osvScanner2026,slsaSpec2026}. \\
Sandbox escape or tool overreach & Least-privilege tools, explicit write scope, generated-zone protections. & Changed-path report and policy decision. \\
Unsafe code and FFI & Unsafe ledger, boundary tests, memory-safety review where the stack exposes unsafe surfaces. & Ledger entry, invariant tests, owner approval. \\
False-green security & Security lane mapped by path and risk, not by agent choice. & CI job link and report artifact. \\
CI downgrade & Pin actions, minimize workflow permissions, require evidence uploads, and diff branch protections. & Workflow proof, permission diff, and release gate receipt. \\
\bottomrule
\end{tabularx}
\end{table*}

Permission policy should be boring. Agents should not receive broad filesystem, network, or credential access merely because a task is urgent. High-risk actions need named lanes and evidence: dependency install, secret material, generated-zone writes, migrations, release artifacts, and browser automation against privileged surfaces. Missing tools are not silently green: doctor diagnostics and security evidence record whether a tool was absent, advisory, failed, or blocking. \verb|--strict| turns policy-blocking security status into a failing lane.
